% =============================================================================
% Odia Phonetic Romanization — Reference Sheet
% Compile with LuaLaTeX (NOT pdflatex — Odia requires Unicode + HarfBuzz shaping):
%     lualatex romanization-reference.tex
% =============================================================================
\documentclass[11pt, a4paper]{article}

\usepackage[margin=1.9cm, headheight=15pt]{geometry}
\usepackage{microtype}
\usepackage{parskip}
\usepackage{booktabs}
\usepackage{tabularx}
\usepackage{array}
\usepackage{enumitem}
\usepackage{mdframed}
\usepackage{xcolor}
\usepackage{titlesec}
\usepackage{fancyhdr}

% --- Fonts: Odia needs HarfBuzz shaping for conjuncts and nukta ---------------
\usepackage{fontspec}
\newfontfamily\odiafont{Oriya Sangam MN}[Renderer=HarfBuzz, Script=Oriya, Scale=1.18]
\newcommand{\od}[1]{{\odiafont #1}}          % inline Odia
\newcommand{\odb}[1]{{\odiafont\large #1}}   % Odia, display size for tables

% --- House palette (MIDSEA accent variant) -----------------------------------
\definecolor{headerblue}{RGB}{0,70,127}
\definecolor{noteblue}{RGB}{220,235,245}
\definecolor{warnyellow}{RGB}{255,248,220}
\definecolor{codegrey}{RGB}{245,245,245}
\definecolor{midseapurple}{RGB}{75,0,130}
\definecolor{midseamagenta}{RGB}{199,21,133}
\definecolor{midseaorange}{RGB}{230,126,0}
\definecolor{midseagreen}{RGB}{0,140,90}
\definecolor{keygreen}{RGB}{223,242,233}

\titleformat{\section}{\large\bfseries\color{midseapurple}}{\thesection}{0.6em}{}[{\color{midseamagenta}\titlerule[1.2pt]}]
\titleformat{\subsection}{\normalsize\bfseries\color{headerblue}}{\thesubsection}{0.5em}{}

\pagestyle{fancy}
\fancyhf{}
\rhead{\small\color{midseapurple}Odia Romanization Reference}
\lhead{\small\color{midseapurple}Odia\_transliteration}
\rfoot{\small Page \thepage}
\lfoot{\small v2 draft --- \texttt{core/rules/symbols.toml} + \texttt{Appv2/rules/*.delta.toml}}
\renewcommand{\headrulewidth}{0.4pt}
\renewcommand{\footrulewidth}{0.2pt}

\newmdenv[backgroundcolor=keygreen, linecolor=midseagreen, linewidth=1pt, roundcorner=4pt,
  innerleftmargin=10pt, innerrightmargin=10pt, innertopmargin=8pt, innerbottommargin=8pt]{keybox}
\newmdenv[backgroundcolor=warnyellow, linecolor=midseaorange, linewidth=1pt, roundcorner=4pt,
  innerleftmargin=10pt, innerrightmargin=10pt, innertopmargin=8pt, innerbottommargin=8pt]{warnbox}
\newmdenv[backgroundcolor=noteblue, linecolor=headerblue, linewidth=1pt, roundcorner=4pt,
  innerleftmargin=10pt, innerrightmargin=10pt, innertopmargin=8pt, innerbottommargin=8pt]{notebox}

\newcommand{\kidea}{{\color{midseagreen}\textbf{\textsf{Key idea.}}}\ }
\newcommand{\kpit}{{\color{midseaorange!85!black}\textbf{\textsf{Pitfall.}}}\ }
\newcommand{\knote}{{\color{headerblue}\textbf{\textsf{Note.}}}\ }

\usepackage{hyperref}
\hypersetup{
  colorlinks = true, linkcolor = midseapurple, urlcolor = headerblue,
  pdftitle = {Odia Phonetic Romanization Reference Sheet},
  pdfauthor = {Amrit Mohapatra}
}

\setlength{\tabcolsep}{5pt}
\renewcommand{\arraystretch}{1.35}

% =============================================================================
\begin{document}
% =============================================================================

\begin{center}
  {\LARGE\bfseries\color{midseapurple} Typing Odia Phonetically}\\[0.3em]
  {\large A reference sheet: what to type in English to get \od{ଓଡ଼ିଆ}}\\[0.4em]
  {\large\color{midseamagenta} v2 --- with Sanskrit in Odia script}\\[0.5em]
  {\small Draft v0.2 \quad$\cdot$\quad \today}
\end{center}

\vspace{0.5em}

\begin{warnbox}
\kpit This sheet is a \textbf{draft scheme}, derived from ITRANS conventions and
published Odia orthography references. It has \textbf{not been reviewed by a native Odia
speaker}. Treat every mapping as a proposal, not a fact.
\end{warnbox}

\begin{notebox}
\knote \textbf{What v2 adds.} Everything in v1 still holds --- this document is a
superset, so nothing you already knew has changed. Sections \ref{sec:sanskrit} onward
cover Sanskrit written in Odia script: the register of stotras, Puranas,
Upanishads and Jagannatha devotional verse. Those additions live in
\texttt{Appv2/rules/*.delta.toml} and never touch \texttt{core/rules/symbols.toml},
which remains the single source of truth for everyday Odia.
\end{notebox}

% =============================================================================
\section{How it works}
% =============================================================================

Odia is an \emph{abugida}, not an alphabet. That one fact explains most of the rules
below.

\begin{keybox}
\kidea Every consonant already carries a built-in vowel. \od{କ} on its own is
\texttt{ka}, not \texttt{k}. To attach a different vowel you add a \textbf{matra}
(a vowel sign). To attach \emph{no} vowel --- so the consonant joins the next one ---
the script uses a \textbf{halant}, which produces a conjunct.
\end{keybox}

There are two typing modes.

\begin{description}[leftmargin=1.4cm, style=nextline, itemsep=0.3em]
  \item[Natural mode \normalfont(default)] Type it the way it sounds. Case is ignored.
  Several spellings reach the same word, and the engine offers a ranked candidate list
  --- you pick. This is what you will use day to day.
  \item[Strict mode \normalfont(\texttt{-{}-strict})] Full ITRANS. Case matters. Exactly
  one deterministic answer, no candidates. Use it when you need a specific character and
  do not want to argue with the ranking.
\end{description}

\subsection*{And, in v2, three \emph{ranking} profiles}

Orthogonal to the two modes above. These do not change what is \emph{reachable} --- every
profile can produce every character --- only what the candidate list offers first.

\begin{description}[leftmargin=1.4cm, style=nextline, itemsep=0.3em]
  \item[Everyday] The original engine, unchanged.
  \item[Odia + Sanskrit \normalfont(default)] Everyday Odia leads; Sanskrit spellings sit
  a slot or two down but are always present. Costs nothing in everyday accuracy.
  \item[Sloka] For verse. A word-final \od{ଂ} beats \od{ମ}, \od{ଳ} beats \od{ଲ}, and the
  Sanskrit vocalics become typeable.
\end{description}

\begin{keybox}
\kidea The profiles only diverge on words \textbf{nothing in the dictionary knows}. For an
attested word the lexicon wins in every profile --- \texttt{balam} is \od{ବଳଂ} in all
three, because \od{ବଳଂ} is a real word and a rule guess should not outrank one.
\end{keybox}

In the tables below, the \textbf{Type} column is natural mode and \textbf{Strict} is the
ITRANS form. Where they are identical, Strict is left blank.

% =============================================================================
\section{Vowels}
% =============================================================================

Each vowel has two shapes: an \textbf{independent} letter, used at the start of a word or
after another vowel, and a \textbf{matra}, used after a consonant. You type the same
thing for both --- the engine picks the right shape from context.

\begin{center}
\begin{tabular}{c c l l l}
\toprule
Independent & On \od{କ} & Type & Strict & Example \\
\midrule
\odb{ଅ}  & \odb{କ}   & \texttt{a}   &                & \od{ଅଛ} \texttt{achha} \\
\odb{ଆ}  & \odb{କା}  & \texttt{aa}  & \texttt{A}     & \od{ଆମ} \texttt{aama} \\
\odb{ଇ}  & \odb{କି}  & \texttt{i}   &                & \od{ଇଏ} \texttt{ie} \\
\odb{ଈ}  & \odb{କୀ}  & \texttt{ii}, \texttt{ee} & \texttt{I} & \od{ନଈ} \texttt{naii} \\
\odb{ଉ}  & \odb{କୁ}  & \texttt{u}   &                & \od{ଉଠ} \texttt{utha} \\
\odb{ଊ}  & \odb{କୂ}  & \texttt{uu}, \texttt{oo} & \texttt{U} & \od{ମୂଳ} \texttt{muula} \\
\odb{ଋ}  & \odb{କୃ}  & \texttt{ru}, \texttt{ri} & \texttt{R\^{}i} & \od{ଋଷି} \texttt{rushi} \\
\odb{ଏ}  & \odb{କେ}  & \texttt{e}   &                & \od{ଏକ} \texttt{eka} \\
\odb{ଐ}  & \odb{କୈ}  & \texttt{ai}, \texttt{oi} &     & \od{ଐକ୍ୟ} \texttt{aikya} \\
\odb{ଓ}  & \odb{କୋ}  & \texttt{o}   &                & \od{ଓଠ} \texttt{otha} \\
\odb{ଔ}  & \odb{କୌ}  & \texttt{au}, \texttt{ou} &     & \od{ଔଷଧ} \texttt{aushadha} \\
\bottomrule
\end{tabular}
\end{center}

\begin{warnbox}
\kpit \textbf{Double the vowel to lengthen it.} This is the single most common reason a
word comes out wrong. \texttt{a} is \od{ଅ} but \texttt{aa} is \od{ଆ}, so \od{ଭାଗ୍ୟ} is
\texttt{bhaagya}, not \texttt{bhagya}. Likewise \texttt{ii}/\texttt{ee} for \od{ଈ} and
\texttt{uu}/\texttt{oo} for \od{ଊ}. Frequent words are recognised either way, but for
anything else the doubling matters.
\end{warnbox}

\begin{notebox}
\knote The rare vowels \od{ୠ}~\texttt{RR}, \od{ଌ}~\texttt{Lu} and \od{ୡ}~\texttt{LL}
exist in the scheme but almost never appear in modern Odia. They are available in strict
mode only.
\end{notebox}

% =============================================================================
\section{Consonants}
% =============================================================================

The five \emph{vargas} (articulation groups). Read each row left to right: unaspirated,
aspirated, voiced, voiced-aspirated, nasal.

\begin{center}
\begin{tabular}{l c c c c c}
\toprule
& \multicolumn{5}{c}{\textit{Position in the mouth}} \\
\cmidrule(l){2-6}
Group & \multicolumn{2}{c}{voiceless} & \multicolumn{2}{c}{voiced} & nasal \\
\midrule
Velar     & \odb{କ} \texttt{ka} & \odb{ଖ} \texttt{kha} & \odb{ଗ} \texttt{ga} & \odb{ଘ} \texttt{gha} & \odb{ଙ} \texttt{nga} \\
Palatal   & \odb{ଚ} \texttt{cha} & \odb{ଛ} \texttt{chha} & \odb{ଜ} \texttt{ja} & \odb{ଝ} \texttt{jha} & \odb{ଞ} \texttt{nya} \\
Retroflex & \odb{ଟ} \texttt{Ta} & \odb{ଠ} \texttt{Tha} & \odb{ଡ} \texttt{Da} & \odb{ଢ} \texttt{Dha} & \odb{ଣ} \texttt{Na} \\
Dental    & \odb{ତ} \texttt{ta} & \odb{ଥ} \texttt{tha} & \odb{ଦ} \texttt{da} & \odb{ଧ} \texttt{dha} & \odb{ନ} \texttt{na} \\
Labial    & \odb{ପ} \texttt{pa} & \odb{ଫ} \texttt{pha} & \odb{ବ} \texttt{ba} & \odb{ଭ} \texttt{bha} & \odb{ମ} \texttt{ma} \\
\bottomrule
\end{tabular}
\end{center}

\begin{warnbox}
\kpit \textbf{Retroflex vs dental} is the one distinction English does not have, and the
main reason you will see multiple candidates. Retroflex (\od{ଟ ଠ ଡ ଢ ଣ}) is made with the
tongue curled back; dental (\od{ତ ଥ ଦ ଧ ନ}) with the tongue on the teeth. In strict mode a
\textbf{capital} letter means retroflex: \texttt{Ta} is \od{ଟ}, \texttt{ta} is \od{ତ}. In
natural mode case is ignored and both are offered --- the dictionary decides which is
more likely.
\end{warnbox}

\subsection*{Everything else}

\begin{center}
\begin{tabular}{c l l c l l}
\toprule
Letter & Type & Notes & Letter & Type & Notes \\
\midrule
\odb{ଯ} & \texttt{ya} & sounds like \emph{ja} & \odb{ଶ} & \texttt{sha} & palatal \\
\odb{ର} & \texttt{ra} &                        & \odb{ଷ} & \texttt{Sha} & retroflex \\
\odb{ଲ} & \texttt{la} &                        & \odb{ସ} & \texttt{sa}  & dental \\
\odb{ଳ} & \texttt{La} & retroflex \emph{l}     & \odb{ହ} & \texttt{ha}  & \\
\odb{ବ} & \texttt{ba}, \texttt{va} &           & \odb{ୱ} & \texttt{wa}  & \od{ସ୍ୱର} \texttt{swara} \\
\bottomrule
\end{tabular}
\end{center}

\subsection*{The three special letters}

These carry a \emph{nukta} (a dot below) or are otherwise distinct. They matter more in
Odia than in most Indic scripts --- the first one appears in the name of the language.

\begin{center}
\begin{tabular}{c l l l}
\toprule
Letter & Type & Sound & Example \\
\midrule
\odb{ଡ଼} & \texttt{rr}, \texttt{d}    & retroflex flap        & \od{ଓଡ଼ିଶା} \texttt{odishaa}, \od{ବଡ଼} \texttt{bada} \\
\odb{ଢ଼} & \texttt{rh}, \texttt{dh}   & aspirated flap        & \od{ପଢ଼ା} \texttt{paRhaa} \\
\odb{ୟ}  & \texttt{y} (not word-initial) & true \emph{y} sound & \od{ରାଜ୍ୟ} \texttt{raajya} \\
\bottomrule
\end{tabular}
\end{center}

\begin{warnbox}
\kpit \od{ଯ} and \od{ୟ} are different letters, and Odia uses them the opposite way round
from Hindi. \od{ଯ} is pronounced like \emph{ja}; \od{ୟ} is the one that actually sounds
like \emph{y}. So typing \texttt{y} gives \od{ଯ} at the \emph{start} of a word
(\texttt{yoga} \(\rightarrow\) \od{ଯୋଗ}) and \od{ୟ} anywhere else
(\texttt{raajya} \(\rightarrow\) \od{ରାଜ୍ୟ}, \texttt{nyaaya} \(\rightarrow\) \od{ନ୍ୟାୟ}).
Both stay in the candidate list, so you can override either way.
\end{warnbox}

% =============================================================================
\section{Signs and marks}
% =============================================================================

These are combining marks, so they are shown here attached to \od{କ} --- on their own
they would be almost invisible.

\begin{center}
\begin{tabular}{c c l l l}
\toprule
Mark & On \od{କ} & Name & Type & Example \\
\midrule
\odb{ଂ} & \odb{କଂ} & anusvara     & \texttt{n}, \texttt{m}, \texttt{M} & \od{ଅଂଶ} \texttt{ansha} \\
\odb{ଁ} & \odb{କଁ} & candrabindu  & \texttt{n\textasciitilde}, \texttt{Mn} & \od{ମୁଁ} \texttt{mun} \\
\odb{ଃ} & \odb{କଃ} & visarga      & \texttt{H}                         & \od{ଦୁଃଖ} \texttt{duHkha} \\
\odb{୍}  & \odb{କ୍} & halant       & automatic                          & see below \\
\bottomrule
\end{tabular}
\end{center}

Anusvara nasalises \emph{before a following consonant}; candrabindu nasalises the vowel
itself; visarga is a breathy release; halant strips the built-in vowel so a conjunct can
form.

% =============================================================================
\section{Conjuncts}
% =============================================================================

When two consonants meet with no vowel between them, they fuse into a single conjunct
glyph. \textbf{You do not type the halant} --- just type the consonants back to back and
the engine forms the conjunct.

\begin{center}
\begin{tabular}{l c l c l c}
\toprule
Type & Get & Type & Get & Type & Get \\
\midrule
\texttt{kta}  & \odb{କ୍ତ} & \texttt{pra}  & \odb{ପ୍ର} & \texttt{nta}  & \odb{ନ୍ତ} \\
\texttt{ksha} & \odb{କ୍ଷ} & \texttt{shra} & \odb{ଶ୍ର} & \texttt{Nda}  & \odb{ଣ୍ଡ} \\
\texttt{jny}  & \odb{ଜ୍ଞ} & \texttt{swa}  & \odb{ସ୍ୱ} & \texttt{mbha} & \odb{ମ୍ଭ} \\
\texttt{tra}  & \odb{ତ୍ର} & \texttt{nka}  & \odb{ଙ୍କ} & \texttt{ddha} & \odb{ଦ୍ଧ} \\
\texttt{sta}  & \odb{ସ୍ତ} & \texttt{ShTa} & \odb{ଷ୍ଟ} & \texttt{lla}  & \odb{ଲ୍ଲ} \\
\texttt{gya}  & \odb{ଗ୍ୟ} & \texttt{nya}  & \odb{ନ୍ୟ} & \texttt{dya}  & \odb{ଦ୍ୟ} \\
\bottomrule
\end{tabular}
\end{center}

\begin{keybox}
\kidea \textbf{A typed \texttt{n} is five different letters}, depending on what
follows it, because a nasal takes the place of articulation of the next consonant:
\texttt{n}+\texttt{k/g} \(\rightarrow\) \od{ଙ} (\od{ଙ୍କ}, \od{ଗଙ୍ଗ});
\texttt{n}+\texttt{ch/j} \(\rightarrow\) \od{ଞ} (\od{ପଞ୍ଚମ});
\texttt{n}+\texttt{T/D} \(\rightarrow\) \od{ଣ} (\od{ମୁଣ୍ଡ});
\texttt{n}+\texttt{t/d} \(\rightarrow\) \od{ନ} (\od{ଅନ୍ତ});
\texttt{n}+\texttt{p/b} \(\rightarrow\) \od{ମ} (\od{ସମ୍ବାଦ}).
The anusvara \od{ଂ} is offered alongside all of them (\od{ସଂସ୍କାର}).
So the honorific is simply \texttt{nka}: \texttt{baabunka} \(\rightarrow\) \od{ବାବୁଙ୍କ}.
\end{keybox}

\begin{warnbox}
\kpit \texttt{ny} and \texttt{gy} are read as \emph{two} letters, not as \od{ଞ} and
\od{ଜ୍ଞ}, because in real words that is nearly always what they are ---
\texttt{dhanyabaada} \(\rightarrow\) \od{ଧନ୍ୟବାଦ}, \texttt{bhaagya} \(\rightarrow\)
\od{ଭାଗ୍ୟ}. For the single letters type \texttt{\textasciitilde n} \(\rightarrow\) \od{ଞ}
and \texttt{jny} \(\rightarrow\) \od{ଜ୍ଞ}.
\end{warnbox}

\begin{notebox}
\knote To force a bare consonant with a visible halant and \emph{no} conjunct --- which
you occasionally want at the end of a word --- append an underscore:
\texttt{kamal\_} gives \od{କମଲ୍} rather than \od{କମଲ}.
\end{notebox}

% =============================================================================
\section{Sanskrit in Odia script}
\label{sec:sanskrit}
% =============================================================================

The Odia script is used to write Sanskrit as well as Odia --- stotras, Puranas,
Upanishads, the Bhagavata, Jayadeva's Gita Govinda. That register needs characters
and spellings that everyday Odia rarely touches, and everything in this section exists
for it.

\subsection*{Four marks you type directly}

These four cannot come from the rule tables at all: a roman key must be lowercase
\texttt{a--z}, and a mark that opens a word is unreachable by construction. So they are
substituted before the engine sees the text.

\begin{center}
\begin{tabular}{l c l l}
\toprule
Type & Get & Name & Example \\
\midrule
\texttt{|}  & \odb{।} & danda        & \texttt{balam |}  $\rightarrow$ \od{ବଳଂ ।} \\
\texttt{||} & \odb{॥} & double danda & \texttt{balam ||} $\rightarrow$ \od{ବଳଂ ॥} \\
\texttt{'}  & \odb{ଽ} & avagraha     & \texttt{so'ham}   $\rightarrow$ \od{ସୋଽହଂ} \\
\texttt{:}  & \odb{ଃ} & visarga      & \texttt{du:kha}   $\rightarrow$ \od{ଦୁଃଖ} \\
\bottomrule
\end{tabular}
\end{center}

\begin{notebox}
\knote \textbf{Two deliberate asymmetries.}

The apostrophe only fires \textbf{after \texttt{o} or \texttt{e}}, so
\texttt{don't} and \texttt{Krishna's} survive untouched. That is not a hedge: Sanskrit
sandhi puts avagraha only after those two vowels, and across every occurrence in an 18.4M
token Odia corpus, 98\% follow \od{ୋ} or \od{େ}.

The colon becomes the \emph{letter} \texttt{h}, not \od{ଃ} directly. Emitting the Odia
character would split the word in two and each half would be transliterated blind ---
\texttt{du:kha} came out \od{ଦୁ}\,\od{ଃ}\,\od{ଖା} even though \od{ଦୁଃଖ} is in the
dictionary. Avagraha has no such problem, because it only ever sits at a seam
\emph{between} two words, where splitting is the correct analysis.
\end{notebox}

\subsection*{The Sanskrit vocalics}

Only \od{ୃ} is in regular use in modern Odia. The rest exist in the script specifically
for Sanskrit transcription, and are reachable by keystroke in Sloka mode or from the
insert palette in any mode.

\begin{center}
\begin{tabular}{c c l l l}
\toprule
Letter & On \od{କ} & Name & Type & Note \\
\midrule
\odb{ଋ} & \odb{କୃ} & vocalic r  & \texttt{ri}, \texttt{ru}  & common in modern Odia \\
\odb{ୠ} & \odb{କୄ} & vocalic rr & \texttt{rri}              & Sloka mode / palette \\
\odb{ଌ} & \odb{କୢ} & vocalic l  & \texttt{lri}              & Sloka mode / palette \\
\odb{ୡ} & \odb{କୣ} & vocalic ll & \texttt{llri}             & Sloka mode / palette \\
\bottomrule
\end{tabular}
\end{center}

\begin{notebox}
\knote These are Sloka-only for a measured reason. The key \texttt{rri} sits on top of
the existing two-character key \texttt{rr} $\rightarrow$ \od{ଡ଼}, which is what made 47
conjuncts reachable in v1. \od{ୠ} occurs 22 times in 527k word types; \od{ଡ଼} is in the
language's own name. Paying candidate-list noise on \od{ଡ଼} everywhere to reach \od{ୠ} is
a bad trade outside verse.
\end{notebox}

\subsection*{Homorganic nasal conjuncts}

% NOTE: \textemdash{} rather than "---" on this line only.
% The "---" ligature silently fails to form when the run sits immediately next
% to an \od{} font switch, and renders as three separate hyphens. Everywhere
% else in this file "---" is fine; it is specifically the Odia-adjacent case.
A nasal before a stop assimilates in the script \textemdash{} \od{ଙ} before velars,
\od{ଞ} before palatals \textemdash{} but nobody types it that way, so a plain
\texttt{n} reaches all of them.
Together these are 3.6\% of all Odia text, and \od{ଙ୍କ} alone is the third most common
conjunct in the language.

\begin{center}
\begin{tabular}{l c l c l c l c}
\toprule
Type & Get & Type & Get & Type & Get & Type & Get \\
\midrule
\texttt{nk}  & \odb{ଙ୍କ} & \texttt{nkh}  & \odb{ଙ୍ଖ} & \texttt{ng} & \odb{ଙ୍ଗ} & \texttt{ngh} & \odb{ଙ୍ଘ} \\
\texttt{nch} & \odb{ଞ୍ଚ} & \texttt{nchh} & \odb{ଞ୍ଛ} & \texttt{nj} & \odb{ଞ୍ଜ} & \texttt{njh} & \odb{ଞ୍ଝ} \\
\bottomrule
\end{tabular}
\end{center}

\subsection*{Subjoined \od{ୱ}}

Sanskrit \emph{va} is written \od{ବ} standing alone --- \od{ଦେବ}, \od{ବିଷ୍ଣୁ} --- but
\od{ୱ} when subjoined under another consonant. This is not a nicety: \od{ସ୍ୱାମୀ}
outnumbers \od{ସ୍ବାମୀ} 3.3 to 1 in real Odia text, and \od{ବିଶ୍ୱ} beats \od{ବିଶ୍ବ} by 7.7
to 1.

\begin{center}
\begin{tabular}{l c l c l c l c}
\toprule
Type & Get & Type & Get & Type & Get & Type & Get \\
\midrule
\texttt{sv}, \texttt{sw} & \odb{ସ୍ୱ} & \texttt{dv}, \texttt{dw} & \odb{ଦ୍ୱ} &
\texttt{tv}, \texttt{tw} & \odb{ତ୍ୱ} & \texttt{shv} & \odb{ଶ୍ୱ} \\
\bottomrule
\end{tabular}
\end{center}

\begin{notebox}
\knote These are spelled out as their own keys because the rule language has no context
meaning ``immediately after a halant''. The available contexts are word-initial,
non-initial, intervocalic and word-final, and all are evaluated against the \emph{roman}
you typed, not against the Odia being assembled. Adding a post-halant context is the one
engine change worth making later.
\end{notebox}

\subsection*{\od{ଳ} against \od{ଲ}, measured}

Sanskrit \emph{la} in an Odia tatsama word is usually \od{ଳ}. This is the single most
visible difference between v1 and v2 output, so it is worth stating what it rests on.
Counting every word type in an 18.4M token corpus, \od{ଳ} against \od{ଲ} for the same
skeleton:

\begin{center}
\begin{tabular}{l r l}
\toprule
Position & Ratio & Prefers \\
\midrule
word-final    & 2.37 & \odb{ଳ} \\
intervocalic  & 1.37 & \odb{ଳ} \\
in a conjunct & 1.21 & \odb{ଳ} \\
subjoined     & 0.23 & \odb{ଲ} \\
word-initial  & 0.05 & \odb{ଲ} \\
\bottomrule
\end{tabular}
\end{center}

Individual words are far more lopsided than the aggregate --- \od{ଜଳ} 7.7$\times$,
\od{ଫଳ} 9.3$\times$, \od{ବଳ} 3.6$\times$ --- but the aggregate is the number the weights
answer to.

\begin{warnbox}
\kpit ``Word-final'' here means what the \emph{engine} means: nothing follows in the
roman. \od{ଥିଲା} is typed \texttt{thilaa}, so its \od{ଲ} is \textbf{intervocalic}, not
word-final. Measuring it the other way sweeps every \od{-ଲା} verb form into the wrong
bucket and drags this ratio from 2.37 down to 0.99 --- which would justify exactly the
opposite weight.
\end{warnbox}

\subsection*{The insert palette}

Some marks are too rare to deserve a keystroke and would only add noise to the candidate
list. Those live on a tap-to-insert strip above the output: the vocalics above,
\od{ଂ} \od{ଁ} \od{ଃ} \od{ଽ}, the halant and nukta, bare \od{ଞ} \od{ଙ} \od{ୱ}, the two
danda, the pranava \od{ଓଁ}, and the Vedic accents U+0951 \emph{udatta} and U+0952
\emph{anudatta} for swara-marked chanting text.

Bare \od{ଞ} is palette-only by design. Giving it a \texttt{ny} key was tried and reverted:
it made the correct \od{ନ୍ୟ} conjunct pay the digraph-split penalty and dropped
\texttt{dhanyabaada} out of the candidate list entirely.

\begin{notebox}
\knote \textbf{Deliberately not included}, so nobody re-adds them. U+0B55 \textsc{oriya
sign overline} reads like a Vedic mark and is not one --- Unicode added it in 13.0 for the
\textbf{Kui} language. U+0B56 and U+0B57 are pieces of a decomposed \od{ୈ}/\od{ୌ}, which
Unicode says to write precomposed. U+0B72--U+0B77 are fraction signs. Odia has no
dedicated OM character; the pranava is \od{ଓ} plus candrabindu.
\end{notebox}

% =============================================================================
\section{Numerals}
% =============================================================================

\begin{center}
\begin{tabular}{l c c c c c c c c c c}
\toprule
Type   & 0 & 1 & 2 & 3 & 4 & 5 & 6 & 7 & 8 & 9 \\
Get    & \odb{୦} & \odb{୧} & \odb{୨} & \odb{୩} & \odb{୪} & \odb{୫} & \odb{୬} & \odb{୭} & \odb{୮} & \odb{୯} \\
\bottomrule
\end{tabular}
\end{center}

Odia numerals are \textbf{off by default} --- modern written Odia overwhelmingly uses
Arabic digits. Turn them on with the \texttt{digits = "odia"} setting.

% =============================================================================
\newpage
\section{Worked examples}
% =============================================================================

\begin{center}
\begin{tabular}{l l l}
\toprule
Type this & Get this & Meaning \\
\midrule
\texttt{namaskaara}   & \odb{ନମସ୍କାର}      & hello \\
\texttt{kemiti achha} & \odb{କେମିତି ଅଛ}     & how are you \\
\texttt{dhanyabaada}  & \odb{ଧନ୍ୟବାଦ}      & thank you \\
\texttt{mun bhala achhi} & \odb{ମୁଁ ଭଲ ଅଛି} & I am well \\
\texttt{odishaa}      & \odb{ଓଡ଼ିଶା}        & Odisha \\
\texttt{odiaa}        & \odb{ଓଡ଼ିଆ}         & Odia (the language) \\
\texttt{bhubaneswara} & \odb{ଭୁବନେଶ୍ୱର}    & Bhubaneswar \\
\texttt{paani}        & \odb{ପାଣି}          & water \\
\texttt{tume}         & \odb{ତୁମେ}          & you \\
\texttt{bhaagya}      & \odb{ଭାଗ୍ୟ}         & fate, luck \\
\texttt{raajya}       & \odb{ରାଜ୍ୟ}         & state \\
\texttt{nyaaya}       & \odb{ନ୍ୟାୟ}         & justice \\
\bottomrule
\end{tabular}
\end{center}

\begin{keybox}
\kidea You are not expected to spell these correctly on the first try. The engine offers
a ranked list, you choose, and it remembers your choice. Accuracy on the words
\emph{you} personally use climbs quickly --- that is the whole design.
\end{keybox}

% =============================================================================
\section{Five things that trip people up}
% =============================================================================

\begin{enumerate}[leftmargin=1.4cm, itemsep=0.45em]
  \item \textbf{Double the vowel to lengthen it.} By far the most common mistake.
        \od{ଭାଗ୍ୟ} is \texttt{bhaagya}; \texttt{bhagya} would be \od{ଭଗ୍ୟ}, which is not a
        word. Same for \texttt{ii}/\texttt{ee} and \texttt{uu}/\texttt{oo}.

  \item \textbf{The built-in vowel is not silent.} \od{ଭଲ} is \texttt{bhala}, not
        \texttt{bhal}. It is also pronounced closer to \emph{o} than \emph{a} --- which is
        why \od{ଓଡ଼ିଶା} is written ``Odisha'' in English, not ``Adisha''.
  \item \textbf{\od{ଡ଼} is not \od{ଡ}.} The flap \od{ଡ଼} is extremely common in Odia
        (\od{ବଡ଼}, \od{ପଢ଼ା}, \od{ଓଡ଼ିଶା}). Internally it is always stored as \od{ଡ} plus a
        separate nukta character, never as a single precomposed character --- they look
        identical on screen but do not match each other in a search.
  \item \textbf{\od{ଳ} is a real letter,} not a typo for \od{ଲ}. Retroflex \emph{l},
        common in Odia: \od{ମୂଳ}, \od{ଫଳ}, \od{କାଳ}. Type \texttt{La}.
  \item \textbf{Case only matters in strict mode.} In natural mode \texttt{Ta} and
        \texttt{ta} are the same input and you disambiguate by picking a candidate.
\end{enumerate}

% =============================================================================
\section{Fixing this sheet}
% =============================================================================

Every mapping above corresponds to an entry in \texttt{core/rules/symbols.toml}. When a
common word comes out wrong, adjust the weight there rather than working around it:

\begin{center}
\begin{minipage}{0.72\textwidth}
\ttfamily\small
\color{midseapurple}[[consonant]]\color{black}\\
roman = "s"\\
odia~~= [\\
\hspace*{2em}\{ text = "\od{ସ}", w = 1.0 \},\\
\hspace*{2em}\{ text = "\od{ଶ}", w = 0.5 \},\\
\hspace*{2em}\{ text = "\od{ଷ}", w = 0.4 \},\\
]
\end{minipage}
\end{center}

Raising a weight promotes that character in the candidate ranking. After any change, run
\texttt{npm run eval} and confirm the accuracy delta against the committed baseline
before keeping it --- a fix for one word often costs you three others.

\subsection*{The Sanskrit mappings are a \emph{delta}}

Everything in section \ref{sec:sanskrit} lives in
\texttt{Appv2/rules/sanskrit.delta.toml} and \texttt{unified.delta.toml}, which are
applied to \texttt{symbols.toml} by a compiler. Nothing in \texttt{core/} is edited. Each
delta block names an id, the weight, and the counts that justify it:

\begin{center}
\begin{minipage}{0.78\textwidth}
\ttfamily\small
\color{midseapurple}[[consonant.replace]]\color{black}\\
id~~~~= "S2"\\
roman = "l"\\
odia~~= [\\
\hspace*{2em}\{ text = "\od{ଲ}", w = 1.0 \},\\
\hspace*{2em}\{ text = "\od{ଳ}", w = 1.15, context = "word\_final" \},\\
\hspace*{2em}\{ text = "\od{ଳ}", w = 1.05, context = "intervocalic" \},\\
\hspace*{2em}\{ text = "\od{ଳ}", w = 0.02, context = "word\_initial" \},\\
]
\end{minipage}
\end{center}

Regenerate and verify with:

\begin{center}
\begin{minipage}{0.78\textwidth}
\ttfamily\small
npx tsx Appv2/scripts/gen-rules.ts\\
npx tsx Appv2/eval/run.ts
\end{minipage}
\end{center}

\begin{keybox}
\kidea Weights are \textbf{not capped at 1.0}. The engine normalises every candidate path
against the best one, so only \emph{ratios} matter --- which is why a word-final \od{ଂ}
can be given \texttt{w~=~1.6} to put it ahead of \od{ମ} outright.
\end{keybox}

The eval run refuses to report a number against a stale rule file, checks that everyday
Odia has not moved, and proves every character in section \ref{sec:sanskrit} is reachable
by \emph{running the engine} rather than by reading the rule tables.

\vfill
\begin{center}
\small\color{gray}
Generated from \texttt{docs/romanization-reference-v2.tex} \quad$\cdot$\quad
compile with \texttt{lualatex} \quad$\cdot$\quad
font: Oriya Sangam MN
\end{center}

\end{document}
